% Generated draft: requires mathematical and source review.
\documentclass[11pt]{article}
\usepackage[margin=25mm]{geometry}
\usepackage{fontspec}
\setmainfont{Latin Modern Roman}
\usepackage{amsmath,amssymb,mathtools,unicode-math}
\setmathfont{Latin Modern Math}
\usepackage{enumitem,needspace,xurl,hyperref}
\setlength{\parindent}{0pt}
\setlength{\parskip}{5pt}
\setlength{\emergencystretch}{3em}
\newcommand{\N}{\mathbb{N}}
\newcommand{\Z}{\mathbb{Z}}
\newcommand{\Q}{\mathbb{Q}}
\newcommand{\R}{\mathbb{R}}
\newcommand{\C}{\mathbb{C}}
\newcommand{\F}{\mathbb{F}}
\newcommand{\A}{\mathbb{A}}
\newcommand{\PP}{\mathbb{P}}
\newcommand{\E}{\mathbb{E}}
\newcommand{\eps}{\varepsilon}
\newcommand{\dd}{\,\mathrm{d}}
\begin{document}
\section*{UnsolvedMath ID: 30004058}
\textbf{Source problem code:} catalogue entry
\subsection{Definitions and mathematical statement}
Let \(f\) denote the modular form in the supplied question, and let \(\tau\) denote its theta operator. For the indices occurring in the source, write
\[
f_0=f,\qquad f_j=\tau(f_{j-1})=\tau^j(f)\quad(1\leq j\leq p).
\]
The source calls \((f_1,\ldots,f_p)\) the theta cycle of \(f\). Write \(\operatorname{Fil}(g)\) for the filtration of a form \(g\), provided that this filtration is defined in the intended modular-form theory. This notation names the source's filtration; its definition and value space are not supplied.

The requested determination is the behavior of filtration under successive applications of \(\tau\), namely the sequence
\[
\operatorname{Fil}(f),\quad
\operatorname{Fil}(\tau(f)),\quad\ldots,\quad
\operatorname{Fil}(\tau^p(f)),
\]
under the hypotheses defining the intended theta cycle. The source does not specify those hypotheses, the modular-form space, the characteristic, the meaning of \(p\), the theta operator, or the filtration. It also does not state a periodicity identity or the treatment of vanishing iterates. Consequently, this is an unresolved formulation of a determination problem, rather than a fully specified universally quantified assertion. No particular filtration convention, numerical weight-change rule, or cycle-closing identity is imposed here.

The supplied background places the question in the context of unitary Shimura varieties. Its dated assessment distinguishes a signature \((n,1)\) case, for which it attributes an explicit theta-cycle filtration formula to de Shalit and Goren, from a broader unitary setting whose general filtration behavior it labels open. The original question does not specify how these settings delimit its intended scope.
\subsection{Short English statement}
Determine how the filtration of a modular form changes through the source's theta cycle \(\tau(f),\ldots,\tau^p(f)\), including its relation to the filtration of \(f\). The underlying modular-form setting, theta operator, filtration, and cycle conditions require clarification before a precise general statement can be given.
\subsection{Sources}
\begin{itemize}
\item Supplied local problem registry, permanent UnsolvedMath ID \(30004058\), title ``Filtration along Modular Theta Cycles,'' exact source problem code \texttt{catalogue entry}; collection: Oberwolfach Reports Open Problems. The registry records proposed year 2019 and supplies no individual proposer.
\item Supplied report citation: \emph{Arithmetic of Shimura Varieties (2020)}, Oberwolfach Reports, DOI 10.4171/owr/2019/2, pp.\ 65--131. \url{https://doi.org/10.4171/owr/2019/2}. This is the supplied report-level provenance; no exact problem page or passage is supplied.
\item Supplied literature reference: de Shalit--Goren, theta cycles. \url{https://doi.org/10.2140/ant.2019.13.1829}.
\item Supplied literature reference: La Porta, generalized theta operators. \url{https://arxiv.org/abs/2209.03717}.
\item Supplied literature reference: Eischen--Mantovan, entire theta operators. \url{https://academic.oup.com/imrn/article/2022/21/16405/6326546}.
\item Literature assessment checked at \texttt{2026-08-22}: the supplied dated curation labels the problem partially solved, attributes an explicit theta-cycle filtration formula in signature \((n,1)\) to de Shalit and Goren, and describes later work as constructing further weight-shifting operators while general filtration behavior across unitary Shimura varieties remains open. This assessment is supplied dated curation derived from the OWR material, not independently verified here; its time-sensitive status and cited claims require independent verification.
\end{itemize}
\end{document}
